\subsection{Definition of the Functional Scope}
\label{subsec:functional-scope}

Before evaluating GEAR and the studied frameworks, we define the functional
scope of the study independently of the current capabilities of GEAR. This
separation prevents the evaluation scope from being restricted to features
already supported by the tool and, consequently, avoids artificially
increasing its measured coverage.

The candidate functional categories were identified through an analysis of
the scientific literature on \gls{llm}-based agents and \gls{mas}s.
Existing studies describe these systems through dimensions such as agent
profiles, goals, reasoning, planning, memory, action, tool use, perception,
environment interaction, communication, coordination, and capability
acquisition \cite{xi2025rise,guo2024large}. \etal{Wu}~\cite{wu2024autogen}
additionally emphasize configurable agents, human involvement, tool use, and
programmable interaction patterns. At the system level, topology,
orchestration logic, memory management, communication protocols, and
error-handling strategies can significantly influence the behaviour and
performance of a \gls{mas} \cite{emde2026maseval}. Finally,
\etal{Dong}~\cite{dong2024agentops} identify traces, logs, metrics, and runtime
artefacts as important elements for observing agentic systems.

These dimensions constitute the initial set of candidate characteristics. We
then selected those that correspond to the objective of this study: evaluating
the configuration, orchestration, and execution capabilities of
general-purpose frameworks for \gls{llm}-based \gls{mas}s. A
characteristic is included when it:

\begin{itemize}
    \item represents an explicit design or configuration decision;
    \item affects the structure, behaviour, or execution of the resulting
    \gls{mas};
    \item can be identified across multiple general-purpose frameworks; and
    \item can be evaluated independently of a particular application domain,
    model-training process, or deployment infrastructure.
\end{itemize}

Table~\ref{tab:functional-scope} reports the resulting functional scope. A
check mark indicates that a category is included in the study, whereas a cross
indicates that the category is scientifically relevant but outside the scope
of the present evaluation.

\begin{table}[t]
    \centering
    \caption{Functional categories identified from the literature and their
    inclusion in the scope of the study.}
    \label{tab:functional-scope}
    \footnotesize
    \begin{tabularx}{\textwidth}{
        @{}
        p{0.17\textwidth}
        >{\raggedright\arraybackslash}X
        >{\centering\arraybackslash}p{0.09\textwidth}
        p{0.29\textwidth}
        @{}
    }
        \toprule
        \textbf{Category}
        & \textbf{Considered characteristics}
        & \textbf{Included}
        & \textbf{Scope rationale} \\
        \midrule

        Agent specification
        & Identity, role, goals, instructions, capabilities, and contextual
        description
        & \included
        & These elements define the responsibilities and expected behaviour
        of each agent. \\

        Model configuration
        & Model, provider, endpoint, generation parameters, timeout, and retry
        parameters
        & \included
        & These configuration decisions can affect agent behaviour and
        execution across frameworks. \\

        Tasks and planning
        & Task descriptions, objectives, inputs, expected outputs,
        decomposition, and dependencies
        & \included
        & Tasks constitute configurable units of work and define how system
        objectives are operationalized. \\

        Tools
        & Tool declarations, tool schemas, association with agents, and tool
        invocation
        & \included
        & Tools determine how agents act on external resources and extend
        their capabilities. \\

        Memory and state
        & Individual memory, shared memory, workflow state, and propagation of
        intermediate results
        & \included
        & Memory and state determine how information is retained and shared
        during execution. \\

        Data and outputs
        & Textual outputs, JSON outputs, typed models, schemas, and output
        validation
        & \included
        & Output contracts determine how information is exchanged among
        tasks, agents, and external components. \\

        Orchestration and topology
        & Sequential and parallel execution, delegation, handoffs,
        fan-out/fan-in, aggregation, nodes, edges, and dependencies
        & \included
        & These characteristics define the structural organization and
        coordination of the system. \\

        Control flow
        & Conditions, branches, loops, routing decisions, and termination
        conditions
        & \included
        & Control-flow mechanisms determine which components are executed and
        under which conditions. \\

        Execution properties
        & Synchronous and asynchronous execution, execution limits, retries,
        error handling, caching, and human intervention
        & \included
        & These properties influence the runtime semantics and robustness of
        the resulting system. \\

        \midrule

        Communication protocols
        & Explicit messages, message types, communication channels,
        conversation protocols, and negotiation mechanisms
        & \excluded
        & Communication is a central multi-agent dimension, but comparing
        protocol semantics requires a separate interaction-level analysis. \\

        Perception and environment
        & Sensor modalities, perception pipelines, environment models, and
        environment-specific actions
        & \excluded
        & The study focuses on orchestrated software agents rather than
        embodied agents or domain-specific perception mechanisms. \\

        Learning and evolution
        & Online learning, self-improvement, capability acquisition, and
        runtime evolution
        & \excluded
        & The study evaluates system configuration and execution, not model
        training or agent-learning processes. \\

        Observability
        & Traces, logs, metrics, monitoring, runtime events, and execution
        analytics
        & \excluded
        & Observability concerns the operational analysis of executions and
        is assessed separately from functional coverage. \\

        Security and governance
        & Identities, permissions, access control, isolation, policy
        enforcement, and audit mechanisms
        & \excluded
        & These mechanisms require a dedicated security model, threat
        analysis, and evaluation protocol. \\

        \bottomrule
    \end{tabularx}
\end{table}

The retained scope therefore covers the declarative and executable
characteristics required to configure agents, assign tasks and tools, manage
memory and data, organize the system, control its execution flow, and define
its runtime properties. These characteristics are sufficiently general to be
examined across the studied frameworks while remaining directly related to
the configuration and execution concerns addressed by this study.

Some important dimensions identified in the literature are deliberately
excluded. In particular, communication is widely recognized as a central
property of \gls{mas}s
\cite{guo2024large,wu2024autogen,emde2026maseval}. However, the present study
does not compare message formats, communication languages, channels,
conversation protocols, or negotiation strategies. Handoffs, task
dependencies, shared state, and fan-out/fan-in structures remain included, but
they are examined as orchestration mechanisms rather than as communication
protocols.

Perception and environment modelling are excluded because the study focuses
on software agents operating through tasks, tools, and workflows rather than
on embodied agents or domain-specific sensing mechanisms. Learning and runtime
evolution are also excluded because the evaluation does not concern model
training, online learning, or agent self-improvement. Security and governance
would require a dedicated threat model and evaluation protocol and are
therefore left for future work.

Observability is handled separately. Although traces, logs, metrics, and
runtime artefacts are important for analysing agentic systems
\cite{dong2024agentops}, they describe how executions are monitored rather
than the functional configuration of a \gls{mas}. They are therefore
assessed in the operational evaluation rather than in the functional coverage
analysis.

